\chapter{先建立全局地图：谁负责理解，谁负责运动}
\section{你要实现的任务}
小鸭子面对一个包含红、黄、蓝三个球的场景。人可以说“踢红球”“把红色的球踢出去”，也可以说“往后走一点”“左转一下”或“现在看到什么”。球可能在当前视野里，也可能在身后。系统需要看图、理解指令、决定一个动作，执行后重新观察，直到任务结束。

你已有的行走和踢球模型，解决的是身体怎么动。上层视觉语言模型解决的是此刻该选择哪个技能、持续多久、是否需要继续观察。这两部分学习的对象和工作频率不同。把已有底层能力保留下来，可以把本轮实验集中在上层决策。

\begin{figure}[htbp]
\centering
\begin{tikzpicture}[>=Latex,node distance=7mm,
  block/.style={draw=accent,fill=pale,rounded corners=1.5mm,align=center,
  text width=139mm,minimum height=13mm,font=\small\sffamily}]
\node[block] (obs) {观测：front / near 相机、姿态、RGB-D 球位置测量\\用户原始指令 + 最近执行历史};
\node[block,below=of obs] (vlm) {上层决策：预训练 / 微调后的 VLM\\选择一个技能、参数，或输出回答 JSON};
\node[block,below=of vlm] (json) {规定接口：type、skill、seconds、target\_color、reply\\格式校验 + 执行条件检查};
\node[block,below=of json] (low) {底层执行：已有 ONNX 行走 / 踢球策略 + 状态反馈\\MuJoCo 物理仿真，随后恢复站立};
\node[block,below=of low] (new) {重新拍摄、记录反馈、进行下一次上层决策\\独立物理验收判断真实结果};
\draw[->,thick,accent] (obs)--(vlm);
\draw[->,thick,accent] (vlm)--(json);
\draw[->,thick,accent] (json)--(low);
\draw[->,thick,accent] (low)--(new);
\draw[->,thick,accent] (new.east)--++(5mm,0)|-(obs.east);
\end{tikzpicture}
\caption{VLM机器人闭环的总体结构。规则老师另行提供示范，原始采集通过Decision接口执行，导出时才转换为统一五字段标签。一次模型输出只负责下一步。}
\end{figure}

\section{五个阶段的输入和输出}
\begin{longtable}{P{24mm}P{43mm}P{43mm}P{36mm}}
\toprule
阶段 & 输入 & 输出 & 执行者 / 状态\\
\midrule
\endhead
基础预训练与后训练 & 大规模文本、图像、视频及监督或反馈数据 & 已具备视觉语言能力的权重 & Qwen团队；我们下载使用\\
规则示范采集 & 随机场景、语言命令、相机和状态 & 观测、规则动作、执行轨迹、验收结果 & 本机MuJoCo；已完成首批\\
本地监督微调 & 经筛选的观测—目标JSON配对 & LoRA适配器、训练状态、指标 & 本机RTX 4050；正在运行\\
模型推理 & 当前观测、指令、允许的历史与状态 & 动作JSON或含reply的回答JSON & 基础VLM已接通；微调LoRA待接入\\
完整任务评测 & 保留的场景和未参与选模的任务 & 成功率、错误类型、耗时、接触证据 & 微调后闭环尚待完成\\
\bottomrule
\end{longtable}

本地微调不需要重新收集互联网级别的大规模图文数据。它从已有权重出发，学习你的技能接口及场景习惯。相反，直接部署一个未经本地微调的 VLM 也有可能输出合法技能，因为它可以阅读提示词；微调的目的是让这种选择更稳定、更贴合实际，而不是首次赋予它读字的能力。

\section{VLM、VLA、规则与技能库各是什么}
\textbf{VLM（视觉语言模型）}接收视觉和语言信息，输出语言或可被程序解释的结构化内容。本项目中的 Qwen3-VL 本体属于 VLM。

\textbf{VLM 驱动的技能调用系统}把这个模型接到预定义动作接口。模型输出 \texttt{move\_forward} 等技能名和参数，执行器再调用底层策略。动作不是模型输出一串自然语言后随意解释，而是按固定接口严格解析。

\textbf{分层视觉语言动作策略}可以用来描述经过动作示范训练、从视觉和语言预测技能的整套系统。是否称为“分层 VLA”取决于使用的术语习惯；关键是写清楚动作层级。本项目输出的是技能层动作，不是每个电机关节的目标角度或力矩。

\textbf{规则老师}用明确程序逻辑生成下一技能。它不通过梯度更新学习，但可以产生训练示范。\textbf{技能库}则给出了动作名称、执行接口、约束和已有控制器。它不是一段“让模型想象鸭子如何运动”的文字。

\begin{notebox}{“固定技能”不等于“固定整段任务”}
左脚踢球的底层动作已经训练好；但是走向哪个球、什么时候转身、如何接近、何时踢球，仍可根据每次新画面改变。你需要学习的主要是这些技能之间的选择与切换。
\end{notebox}

\section{训练、推理和物理控制是三个不同的时钟}
训练中的一次更新，在当前机器上约需22秒，处理四条已经保存的样本。此时不要求真实机器人同步运动，主要工作是计算损失和梯度。

推理中的一次决策，消耗的是看当前图像并生成输出的时间。它可能随着模型、后端、输入长度和显卡状态改变。训练一次更新的22秒不能直接当作部署时一次技能决策的延迟。

底层控制以较高频率读取身体状态并给出关节动作。当前策略控制约50Hz、物理仿真步频200Hz。上层模型等待时仍须维持站立；不能等每个语言token出来才给电机一次指令。

\section{“模型知道什么”的四个来源}
\begin{enumerate}
\item \textbf{预训练权重}：学到图像、词语、颜色、空间关系和语言表达等统计规律；不保证天然理解这只鸭子的精确踢球坐标。
\item \textbf{本次提示词}：列出九个技能、动作参数、坐标方向和输出五个字段。模型通过阅读上下文获知当前接口。
\item \textbf{当前观测和指令}：用户说“红球”，两路相机和辅助测量提供可见球的信息。目标颜色来自原始指令及对话上下文。
\item \textbf{本地示范训练}：把典型输入与合适的技能选择关联起来，更新LoRA或未来的技能预测头。
\end{enumerate}
这也解释了“没有微调，怎么知道技能库”：知道接口可以依靠提示词；稳定掌握接口在复杂场景下的正确使用，需要实际数据和验证。格式由程序设计者规定，模型学习或遵循这个协议。

\section{三个“成功”必须分别记录}
\begin{longtable}{P{31mm}P{54mm}P{66mm}}
\toprule
成功的层次 & 能说明什么 & 不能直接说明什么\\
\midrule
\endhead
训练程序成功运行 & 前向、反向、优化器更新和检查点保存正常 & 小鸭子已经会踢指定球\\
离线输出匹配示范 & 在给定画面下生成的格式、技能或参数符合标签 & 自己行动后仍能完成完整任务\\
闭环物理任务成功 & 从起点出发，实际搜索、接近、踢中目标并站稳 & 任意新场景、真实机器人和所有新说法都能成功\\
\bottomrule
\end{longtable}
后面每个章节都会保留这三个层次。特别是规则老师的277/300随机踢球成功率，是采集质量证据，不能写成微调后VLM的成功率。

\section{本轮实验的边界}
目前覆盖的是固定三色球、已有九种技能、有限范围随机布局，以及相应移动和视觉问答。拾布、张嘴发声、真实语音识别、麦克风和实体机器人部署，不属于这次已完成的数据与验证范围。“现在能看到什么”的示范主要描述本场景中的球，也不能据此宣称已验证开放世界视觉问答。

这是一轮可复现的起点。保留失败记录、区分训练与测试、说明哪些输入来自辅助测量，比用一个好看的损失数值宣称全部学会更有助于判断下一步该改数据、改提示词，还是改模型。
